% ===================================================================
%  TABLEAU N° 5 — La maison et sa construction.
%
%  Traduction, faite en 2026, en francais standard du Quebec, sur
%  l'ido de texto/io. Regles de la colonne : voir 10-tableau-01.tex.
%  Le tableau est un DIALOGUE et l'auteur n'y numerote aucun alinea :
%  toutes les cles portent donc le meme « 01 », chaque replique
%  formant un rang.
%
%  LES HUIT CHOIX PROPRES A CE TABLEAU :
%    (3)   cave ................. sous-sol
%    (20)  cabinets d'aisances .. toilettes
%    (j)   cabinet de toilette .. salle d'eau
%    (1)   salle de bains ....... salle de bain
%    (e)   office ............... garde-manger
%    (f)   cabinet de travail ... bureau
%    (133) valet d'écurie ....... palefrenier
%    (16)  parquets ............. planchers
%
%  DEUX DE CES HUIT SE TIENNENT PAR LA MAIN. « Cave » et « sous-sol »
%  ne nomment pas la meme chose au Quebec : la cave y est le local
%  frais ou l'on garde le vin et les legumes, et l'etage enterre d'une
%  maison s'appelle le sous-sol. Or l'ido definit l'objet lui-meme —
%  « l'espace compris entre les fondations » — et lui donne une porte,
%  un escalier et un soupirail : c'est un sous-sol.
%  De meme « garde-manger » : au Quebec ce n'est pas l'armoire, c'est
%  la piece attenante a la cuisine, exactement ce que l'ido appelle
%  « servisteyo » et le fac-simile « office ».
%
%  QUATRE MOTS DU FAC-SIMILE RESTENT, ET C'EST UN RESULTAT. Corridor,
%  perron, remise, buanderie : le francais de France leur prefererait
%  volontiers couloir, seuil, hangar, lingerie ; le Quebec dit les
%  quatre premiers, en 2026 comme en 1926. On ne remplace pas un mot
%  pour la seule raison qu'il est ancien.
%
%  « PREMIER ETAGE » RESTE AUSSI, et c'est le cas le plus discutable
%  de la colonne. L'usage parle du Quebec compte souvent a
%  l'americaine — le rez-de-chaussee est « le premier », l'etage
%  au-dessus « le deuxieme ». Mais l'ido ecrit « teretajo » puis
%  « unesma etajo », le plan de la planche porte ces deux mots, et
%  l'Office quebecois de la langue francaise recommande la
%  numerotation francaise. On la garde : ici le mot renvoie a une
%  legende gravee, et la legende ne se renumerote pas.
% ===================================================================

%%K t05-apar-1 apar
\VUsaut{6.0mm}
\VUtitre{15.0pt}{60}{TABLEAU N\textsuperscript{o} 5}
\VUsaut{1.4mm}
\VUtitre{11.4pt}{0}{La maison et sa construction.}
\VUsaut{2.2mm}

%%K t05-tit-1 sub
\VUpk{10.2pt}{40}{Bonne rencontre.}
\VUsaut{3.0mm}

%%K t05-01-1 p
\textsc{Jean}. --- Mon oncle, j'aimerais beaucoup savoir comment on construit
une \VUgras{maison}.

%%K t05-01-2 p
\textsc{L'Oncle} \textsuperscript{(80)}. --- C'est bien facile, mon ami; viens
avec moi, je te montrerai celle que monsieur Bernard \textsuperscript{(79)}
fait construire et que j'habiterai.

%%K t05-01-3 p
\textsc{Jean}. --- Et qui donc a dessiné le \VUgras{plan} \textsuperscript{(76)}
de cette maison?

%%K t05-01-4 p
\textsc{L'Oncle}. --- C'est l'\VUgras{architecte} \textsuperscript{(75)}; il a
présenté un dessin très clair, qui a été approuvé, et
l'\VUgras{entrepreneur} \textsuperscript{(77)} a commencé les travaux.

%%K t05-01-5 p
\textsc{Jean}. --- Et qui est l'entrepreneur?

%%K t05-01-6 p
\textsc{L'Oncle}. --- C'est monsieur Leblanc, le père de ton ami Jacques. Il
dirige les ouvriers avec monsieur Leroux, l'architecte.

%%K t05-01-7 p
Tiens! Voici justement monsieur Leblanc avec sa \VUgras{règle}
\textsuperscript{(78)}, et monsieur Leroux avec son plan.

%%K t05-01-8 p
Bonjour, messieurs. Comment allez-vous?

%%K t05-01-9 p
\textsc{Monsieur Leroux}. --- Très bien, merci. Bonjour, Jean; comme cet enfant
a grandi!

%%K t05-01-10 p
\textsc{L'Oncle}. --- Oui, vraiment, c'est un petit homme. Il est très sérieux,
il faut le renseigner sur tout ce qu'il voit. Aujourd'hui, il voudrait savoir
comment on construit une maison.

%%K t05-tit-2 sub
\VUpk{10.2pt}{40}{Les ouvriers.}
\VUsaut{3.0mm}

%%K t05-01-11 p
\textsc{Monsieur Leblanc}. --- Eh bien, viens avec moi, mon petit ami, je vais
te l'expliquer tout de suite, parce que j'aime beaucoup les enfants qui veulent
s'instruire.

%%K t05-01-12 p
Tu vois cet ouvrier qui tient une \VUgras{pelle} \textsuperscript{(82)} à la
main : c'est un \VUgras{terrassier} \textsuperscript{(81)}. Il creuse une
\VUgras{tranchée} \textsuperscript{(46)} pour poser la tuyauterie d'eau. Il a
aussi creusé le sol à l'endroit où se trouve la maison, pour que le maçon puisse
élever les quatre \VUgras{murs}. La partie de ces murs qui est dans la terre
s'appelle les \VUgras{fondations} \textsuperscript{(2)}. L'espace compris entre
les fondations forme le \VUgras{sous-sol} \textsuperscript{(3)}. Ce sous-sol a
une petite fenêtre appelée \VUgras{soupirail} \textsuperscript{(5)}, une
\VUgras{porte} \textsuperscript{(4)} et un escalier.

%%K t05-01-13 p
Le \VUgras{maçon} \textsuperscript{(108)} a continué à construire les murs en
laissant des ouvertures pour les \VUgras{portes} \textsuperscript{(8)} et les
\VUgras{fenêtres} \textsuperscript{(17)}.

%%K t05-01-14 p
\textsc{Jean}. --- Mais ce maçon, je ne le vois pas!

%%K t05-01-15 p
\textsc{Monsieur Leblanc}. --- Il a fini de travailler à la maison. Il est près
de l'\VUgras{écurie} \textsuperscript{(54)}, sur son \VUgras{échafaudage}
\textsuperscript{(110)}; il construit le mur de la \VUgras{buanderie}
\textsuperscript{(56)}.

%%K t05-01-16 p
Il a une truelle, une auge et un \VUgras{fil à plomb} \textsuperscript{(109)}.
Mais tu ne retiendras pas ces noms-là.

%%K t05-01-17 p
\textsc{Jean}. --- Mais oui, parce que je vais tout de suite demander à mon
oncle une petite truelle et une auge, et je fabriquerai moi-même un fil à plomb
avec une ficelle.

%%K t05-tit-3 sub
\VUsaut{3.0mm}
\VUpk{10.2pt}{40}{Les matériaux.}
\VUsaut{3.0mm}

%%K t05-01-18 p
\textsc{L'Oncle}. --- Ah! très bien. Mais dis-moi, Jean, de quoi a-t-on besoin
pour construire une maison?

%%K t05-01-19 p
\textsc{Jean}. --- Il faut des \VUgras{pierres de taille} \textsuperscript{(59)}
et du \VUgras{mortier} \textsuperscript{(65)}.

%%K t05-01-20 p
\textsc{L'Oncle}. --- Exactement. Regarde cet homme au grand chapeau, sur sa
\VUgras{charrette} \textsuperscript{(136)}. C'est un \VUgras{charretier}
\textsuperscript{(135)}. Chaque jour, il apporte ici des \VUgras{blocs de
pierre} \textsuperscript{(60)}. Il décharge sa voiture dans la cour, et les
\VUgras{tailleurs de pierre} \textsuperscript{(83)} taillent les blocs et les
scient avec leur \VUgras{scie à pierre} \textsuperscript{(85)}. Un
\VUgras{manœuvre} \textsuperscript{(146)} les porte au maçon, qui les met en
place.

%%K t05-01-21 p
Pendant ce temps, le \VUgras{gâche-mortier} \textsuperscript{(87)} prépare le
mortier. Il a besoin de \VUgras{sable} \textsuperscript{(62)}, de
\VUgras{chaux} \textsuperscript{(63)} et d'\VUgras{eau} \textsuperscript{(64)}.
La chaux est près de lui dans un \VUgras{sac} \textsuperscript{(71)}; un
\VUgras{aide-maçon} \textsuperscript{(111)} lui apporte le sable dans une
\VUgras{brouette} \textsuperscript{(112)} et l'\VUgras{apprenti}
\textsuperscript{(113)} est à la pompe \textsuperscript{(58)}. Il remplit un
\VUgras{seau} \textsuperscript{(114)}. Le mortier sera bientôt prêt. Le
\VUgras{porteur de mortier} \textsuperscript{(107)} le prend et le monte par
l'échelle sur l'échafaudage, où travaille un maçon qui termine ce côté de la
maison.

%%K t05-01-22 p
Et maintenant, comment s'appelle l'ouvrier qui fait le \VUgras{toit}
\textsuperscript{(31)}?

%%K t05-01-23 p
\textsc{Jean}. --- C'est le \VUgras{couvreur} \textsuperscript{(118)}.

%%K t05-tit-4 sub
\VUsaut{3.0mm}
\VUpk{10.2pt}{40}{Le toit de la maison.}
\VUsaut{3.0mm}

%%K t05-01-24 p
\textsc{Monsieur Leblanc}. --- Oui, mais il y a d'abord le
\VUgras{charpentier} \textsuperscript{(115)}.

%%K t05-01-25 p
Le charpentier fait la \VUgras{charpente} \textsuperscript{(35)}. Il assemble
les \VUgras{poutres} \textsuperscript{(37)} avec des \VUgras{chevilles}
\textsuperscript{(117)}. Ces ouvriers, là-haut, avec leurs larges culottes de
velours, sont des charpentiers. L'un tient une petite poutre et l'autre coupe
une cheville avec sa \VUgras{hache} \textsuperscript{(116)}. Celui qui est près
de la \VUgras{cheminée} \textsuperscript{(41)} est le couvreur. Il cloue des
lattes sur les (*) solives, et des \VUgras{ardoises} \textsuperscript{(66)} sur
les lattes. Parfois, il pose des tuiles.

%%K t05-noto-1 noto
\VUnotes{143.2mm}{%
(*) \VUgras{Balk-o} = all. \textit{Balken}, angl. \textit{joist}, fr.
\textit{solive}, it. \textit{travicello}, russe \textit{balka}, esp. et port.
\textit{viga}.
Racine technique jusqu'ici non adoptée, mais qu'il y aurait lieu de proposer
pour rendre le sens du mot français \textit{solive}, qui n'est ni une poutre
(fr. \textit{poutre}) ni une poutrelle. C'est une pièce de bois équarrie qui
soutient un plancher et un plafond.%
}

%%K t05-01-26 p
Le toit de l'écurie est en \VUgras{tuiles} \textsuperscript{(55)}, celui de la
maison est en \VUgras{ardoises} \textsuperscript{(32)} et celui de la véranda
est en \VUgras{zinc} \textsuperscript{(24)}.

%%K t05-01-27 p
\textsc{Jean}. --- Est-ce que le couvreur fait aussi les toits de zinc?

%%K t05-01-28 p
\textsc{Monsieur Leblanc}. --- Non, c'est le \VUgras{zingueur}
\textsuperscript{(121)} ou le \VUgras{plombier}, cet ouvrier qui pose une
\VUgras{lucarne} \textsuperscript{(38)} sur le toit de la maison. Il pose aussi
les \VUgras{chéneaux} \textsuperscript{(43)} et les \VUgras{gouttières}
\textsuperscript{(42)}. Il soude ensemble le zinc et le fer-blanc. C'est lui qui
a fixé le \VUgras{paratonnerre} \textsuperscript{(40)} et la
\VUgras{girouette} \textsuperscript{(39)} sur le \VUgras{pignon}
\textsuperscript{(36)}.

%%K t05-01-29 p
\textsc{Jean}. --- La maison est donc terminée?

%%K t05-tit-5 sub
\VUsaut{3.0mm}
\VUpk{10.2pt}{40}{Les outils.}
\VUsaut{3.0mm}

%%K t05-01-30 p
\textsc{Monsieur Leblanc}. --- Pas tout à fait. Le \VUgras{plâtrier}
\textsuperscript{(99)} construit avec des \VUgras{briques}
\textsuperscript{(61)} les cloisons qui la divisent en différentes pièces. Il
enduit de \VUgras{plâtre} \textsuperscript{(70)} les \VUgras{plafonds}
\textsuperscript{(26)} et les murs. En ce moment, il fait le plafond de la
\VUgras{véranda} \textsuperscript{(33)}. Il a, comme le maçon, une
\VUgras{truelle} \textsuperscript{(100)} et une \VUgras{auge}
\textsuperscript{(101)}.

%%K t05-01-31 p
Ensuite, le menuisier fait les portes, les fenêtres, les \VUgras{planchers}
\textsuperscript{(16)} et les \VUgras{volets} \textsuperscript{(19)}.

%%K t05-01-32 p
Il travaille en ce moment dans la cour, à son \VUgras{établi}
\textsuperscript{(90)}. Il a une \VUgras{scie} \textsuperscript{(94)} pour
scier le \VUgras{bois} \textsuperscript{(69)}, un \VUgras{rabot}
\textsuperscript{(91)}, un \VUgras{valet} \textsuperscript{(92)}, un
\VUgras{ciseau} \textsuperscript{(94 \textit{bis})}, un \VUgras{compas}
\textsuperscript{(95)} et un \VUgras{maillet de bois} \textsuperscript{(95
\textit{bis})}.

%%K t05-01-33 p
\textsc{Jean}. --- Ah! mais c'est plus amusant de faire de la menuiserie que de
la maçonnerie. Mon oncle, tu m'achèteras un établi, une scie et un rabot, parce
que je vais bientôt faire une niche pour Médor.

%%K t05-01-34 p
\textsc{L'Oncle}. --- Oui, mon petit.

%%K t05-01-35 p
\textsc{Jean}. --- Comme je suis content! Et qui est celui qui se tient près de
la porte d'entrée?

%%K t05-tit-6 sub
\VUsaut{3.0mm}
\VUpk{10.2pt}{40}{La peinture.}
\VUsaut{3.0mm}

%%K t05-01-36 p
\textsc{Monsieur Leblanc}. --- C'est un \VUgras{peintre}
\textsuperscript{(102)}. Il est sur son échelle avec son pot de
\VUgras{peinture} \textsuperscript{(103)} et son \VUgras{pinceau}
\textsuperscript{(104)}. Il peint le bois de la porte d'entrée pour qu'il se
conserve plus longtemps.

%%K t05-01-37 p
C'est lui aussi qui pose le papier peint dans les appartements.

%%K t05-01-38 p
Le \VUgras{tapissier} \textsuperscript{(107)}, qui est sur une
\VUgras{échelle} \textsuperscript{(106)}, sur le \VUgras{balcon}
\textsuperscript{(28)}, pose un \VUgras{store} \textsuperscript{(29)}.

%%K t05-01-39 p
L'ouvrier qui se tient près de la grande fenêtre est le \VUgras{vitrier}
\textsuperscript{(96)}. Il fixe les \VUgras{vitres} \textsuperscript{(18)} avec
du \VUgras{mastic} \textsuperscript{(73)}. Cet autre ouvrier, à genoux près de
la porte du sous-sol, est un \VUgras{serrurier} \textsuperscript{(97)}. Il pose
une \VUgras{serrure} \textsuperscript{(6)}. Sa \VUgras{lime}
\textsuperscript{(98)} est près de lui.

%%K t05-01-40 p
\textsc{Jean}. --- Et celui qui frappe avec son \VUgras{marteau}
\textsuperscript{(84)} sur du fer, est-ce aussi un serrurier?

%%K t05-01-41 p
\textsc{Monsieur Leblanc}. --- C'est plutôt un forgeron \textsuperscript{(125)}.
Il forge des \VUgras{pièces de fer} \textsuperscript{(67)} pour
l'\VUgras{électricien} \textsuperscript{(124)}, qui est sur le toit de la
\VUgras{remise} \textsuperscript{(51)}. Il a une \VUgras{forge}
\textsuperscript{(126)}, une \VUgras{enclume} \textsuperscript{(127)} et une
\VUgras{pince} \textsuperscript{(128)}.

%%K t05-01-42 p
L'électricien pose le \VUgras{fil de cuivre} \textsuperscript{(44)} du
téléphone.

%%K t05-tit-7 sub
\VUpk{10.2pt}{40}{Le jardinage.}
\VUsaut{3.0mm}

%%K t05-01-43 p
\textsc{L'Oncle}. --- Au fond de la \VUgras{cour} \textsuperscript{(45)}, près
de la \VUgras{grille} \textsuperscript{(47)} et de la \VUgras{porte cochère}
\textsuperscript{(48)}, tu vois le \VUgras{jardinier} \textsuperscript{(129)}.
Il prépare le petit \VUgras{jardin} \textsuperscript{(49)}. Il a bêché la terre
avec sa \VUgras{bêche} \textsuperscript{(131)}, il sème des graines et ratisse
avec le \VUgras{râteau} \textsuperscript{(130)}. Ensuite, avec son
\VUgras{arrosoir} \textsuperscript{(132)}, il arrosera les
\VUgras{plates-bandes} \textsuperscript{(150)} et les plantes pousseront.

%%K t05-01-44 p
Le \VUgras{palefrenier} \textsuperscript{(133)}, qui vient de soigner le cheval
et de lui donner à manger, nettoie la \VUgras{voiture} \textsuperscript{(52)}
avec une éponge, une \VUgras{pompe à main} \textsuperscript{(134)} et un seau
d'eau. Ensuite, il la rentrera dans la remise et fermera la porte avec le gros
\VUgras{verrou} \textsuperscript{(53)}.

%%K t05-01-45 p
Te voilà content, je pense; tu as eu assez d'explications.

%%K t05-01-46 p
\textsc{Jean}. --- Oui, mon oncle, mais j'aimerais bien voir l'intérieur de la
maison.

%%K t05-01-47 p
\textsc{L'Oncle}. --- Ce sera pour demain. Monsieur Leroux t'expliquera le plan
et te dira le nom de chaque pièce.

%%K t05-tit-8 sub
\VUsaut{3.0mm}
\VUpk{10.2pt}{40}{La distribution intérieure de la maison.}
\VUsaut{2.4mm}

%%K t05-apar-2 apar
{\centering\textit{(Voir le plan.)}\par}
\VUsaut{2.4mm}

%%K t05-01-48 p
\textsc{L'Architecte}. --- Viens, Jean, je vais te faire visiter les
différentes parties de la maison.

%%K t05-01-49 p
Entrons au \VUgras{rez-de-chaussée} \textsuperscript{(7)}. Voilà le bouton de
la \VUgras{sonnette} \textsuperscript{(11)}. Nous montons les trois
\VUgras{marches} \textsuperscript{(14)} du \VUgras{perron}
\textsuperscript{(9)}, nous passons le \VUgras{seuil} \textsuperscript{(10)} de
la \VUgras{porte d'entrée} \textsuperscript{(8)}, et nous voici au pied de
l'\VUgras{escalier} \textsuperscript{(13)}.

%%K t05-01-50 p
\textsc{Jean}. --- C'est le corridor.

%%K t05-01-51 p
\textsc{L'Architecte}. --- Non, c'est le \VUgras{vestibule}
\textsuperscript{(12)}. Le \VUgras{corridor} \textsuperscript{(a)} se trouve au
bout. À droite, voici le \VUgras{salon} \textsuperscript{(b)} et la
\VUgras{salle à manger} \textsuperscript{(c)}; en face de la salle à manger se
trouvent la \VUgras{cuisine} \textsuperscript{(d)} et le
\VUgras{garde-manger} \textsuperscript{(e)}, puis sur le devant le
\VUgras{bureau} \textsuperscript{(f)}. La \VUgras{véranda}
\textsuperscript{(23)}, avec sa \VUgras{porte vitrée} \textsuperscript{(25)},
est à côté du salon.

%%K t05-01-52 p
Nous allons monter l'escalier. Tiens-toi à la \VUgras{rampe}
\textsuperscript{(15)} et compte les marches. Il y en a quatorze. Voici le
\VUgras{palier} \textsuperscript{(m)} : nous sommes au \VUgras{premier étage}
\textsuperscript{(27)}.

%%K t05-tit-9 sub
\VUsaut{3.0mm}
\VUpk{10.2pt}{40}{Le premier étage.}
\VUsaut{3.0mm}

%%K t05-01-53 p
En face de l'escalier, il y a les \VUgras{toilettes} \textsuperscript{(20)} et
la \VUgras{salle d'eau} \textsuperscript{(j)}. À droite, une belle
\VUgras{chambre à coucher} \textsuperscript{(g)} avec deux fenêtres sur la
\VUgras{façade} \textsuperscript{(1)} de la maison. À gauche, c'est la
\VUgras{salle de bain} \textsuperscript{(1)} et la \VUgras{chambre d'amis}
\textsuperscript{(i)} avec une grande \VUgras{alcôve} \textsuperscript{(h)}. À
côté de la chambre à coucher, il y a la \VUgras{chambre des enfants}
\textsuperscript{(k)}. Elle donne sur une vaste \VUgras{terrasse}
\textsuperscript{(21)}. Sous cette terrasse se trouvent d'autres toilettes
\textsuperscript{(20)} et, contre le mur, voici la \VUgras{cloche}
\textsuperscript{(22)} qui sonne le souper.

%%K t05-01-54 p
\textsc{Jean}. --- Montons au \VUgras{deuxième étage}.

%%K t05-01-55 p
\textsc{L'Architecte}. --- Il n'y a pas de deuxième étage. Sous le toit, il n'y
a que des \VUgras{mansardes} \textsuperscript{(33)} et le grenier
\textsuperscript{(34)}. Nous avons fini la visite, nous pouvons redescendre.

%%K t05-01-56 p
\textsc{Jean}. --- Merci, monsieur, c'est très intéressant. Je vais tout de
suite raconter à mon père tout ce que j'ai vu.
\VUsaut{4.0mm}
\VUfilet{22mm}
